\section{GIỚI THIỆU ĐỀ TÀI}

\subsection{Bối cảnh và lý do lựa chọn}

Trong trí tuệ nhân tạo (Artificial Intelligence -- AI), ra quyết định tuần tự là lớp bài toán trong đó hệ thống phải lựa chọn hành động dựa trên trạng thái hiện tại, đồng thời xem xét ảnh hưởng của quyết định đó đối với các trạng thái và kết quả về sau. Học tăng cường (Reinforcement Learning -- RL) là một hướng tiếp cận cho lớp bài toán này. Thay vì được cung cấp sẵn hành động đúng cho từng tình huống, tác tử (agent) tương tác với môi trường, nhận phần thưởng và điều chỉnh cách lựa chọn hành động nhằm tối đa hóa phần thưởng tích lũy \cite{SuttonBarto2018}.

Trên cơ sở đó, nhóm lựa chọn Pacman làm môi trường xây dựng và đánh giá các thuật toán học tăng cường \cite{BerkeleyPacman2024}. Luật chơi trực quan nhưng vẫn thể hiện rõ đặc điểm của một bài toán ra quyết định tuần tự. Tại mỗi bước, tác tử phải lựa chọn hướng di chuyển dựa trên vị trí hiện tại, các vật phẩm còn lại và sự xuất hiện của các mối nguy hiểm trong mê cung. Thay vì sử dụng nguyên trạng một môi trường có sẵn, nhóm trực tiếp thiết kế và cài đặt Mini Pacman với bản đồ, không gian trạng thái, không gian hành động, cơ chế phần thưởng và điều kiện kết thúc riêng cho đề tài. Nhờ đó, nhóm có thể kiểm soát các yếu tố ảnh hưởng đến quá trình học và theo dõi trực tiếp hành vi của tác tử.

Trên môi trường này, nhóm cài đặt Q-Learning, DQN và Double DQN trong cùng một ứng dụng. Q-Learning biểu diễn giá trị của các cặp trạng thái--hành động bằng Q-table \cite{WatkinsDayan1992}. DQN sử dụng mạng nơ-ron để xấp xỉ hàm giá trị hành động, kết hợp replay buffer và target network trong quá trình huấn luyện \cite{Mnih2013,Mnih2015}. Double DQN thay đổi cách xây dựng giá trị mục tiêu nhằm giảm xu hướng đánh giá quá cao (overestimation) của DQN \cite{VanHasselt2016}. Việc triển khai ba thuật toán giúp nhóm khảo sát sự khác biệt giữa phương pháp dạng bảng, phương pháp xấp xỉ bằng mạng nơ-ron và một cải tiến trực tiếp của DQN. Trong quá trình thực hiện đề tài, mã nguồn và các cấu hình được nhóm phát triển, quản lý trong repository DRL-Pacman trên GitHub. Kết quả thực nghiệm cũng được lưu trữ và công bố tại đây \cite{DRLPacmanGitHub}.

\subsection{Mục tiêu nghiên cứu}

Mục tiêu tổng quát của đề tài là xây dựng môi trường Mini Pacman, cài đặt và so sánh Q-Learning, DQN và Double DQN trong bài toán điều khiển tác tử chơi Pacman. Qua đó, nhóm xem xét mối liên hệ giữa nguyên lý của từng thuật toán, cách triển khai trong ứng dụng và kết quả quan sát được từ thực nghiệm.

\noindent Các mục tiêu cụ thể gồm:

\begin{itemize}
    \item Xây dựng môi trường Mini Pacman với không gian trạng thái, không gian hành động, cơ chế phần thưởng và điều kiện kết thúc phù hợp cho quá trình huấn luyện tác tử.
    \item Cài đặt Q-Learning, DQN và Double DQN theo các nguyên lý được trình bày trong những công trình nền tảng, đồng thời tổ chức cấu hình, huấn luyện, đánh giá và lưu mô hình cho từng thuật toán \cite{WatkinsDayan1992,Mnih2013,Mnih2015,VanHasselt2016}.
    \item Thực hiện chín cấu hình thí nghiệm, tương ứng với ba thuật toán và ba mức learning rate, trong cùng môi trường và quy trình đánh giá.
    \item Thu thập và so sánh kết quả bằng các chỉ số reward, completion rate, win rate và steps. Loss chỉ được sử dụng khi phân tích DQN và Double DQN.
    \item Xây dựng demo các tác tử sau huấn luyện, từ đó kết hợp quan sát định tính với kết quả định lượng để đánh giá ưu điểm, hạn chế của từng phương pháp và giới hạn của thực nghiệm.
\end{itemize}

\subsection{Đối tượng và phạm vi nghiên cứu}

Đối tượng nghiên cứu của đề tài là ba thuật toán RL dựa trên giá trị: Q-Learning dạng bảng, DQN và Double DQN. Báo cáo tập trung vào nguyên lý hoạt động, cách cài đặt, quá trình huấn luyện và kết quả thực nghiệm của ba thuật toán trong ứng dụng do nhóm xây dựng.

Thực nghiệm được tiến hành trên bản đồ \texttt{medium} có kích thước $15 \times 15$ ô của môi trường Mini Pacman. Tác tử Pacman lựa chọn một trong bốn hướng di chuyển và tương tác với thức ăn, viên sức mạnh, trái cây cùng ba con ma trong mê cung. Q-Learning sử dụng trạng thái rời rạc làm khóa của Q-table, trong khi DQN và Double DQN nhận vector đặc trưng được tạo từ trạng thái môi trường. Đề tài không nghiên cứu phương pháp học trực tiếp từ ảnh chụp màn hình.

Phạm vi thực nghiệm gồm chín cấu hình được tạo từ ba thuật toán và ba mức learning rate lựa chọn trước. Mỗi cấu hình được huấn luyện với một seed. Tại mỗi thời điểm đánh giá định kỳ, mô hình được chạy trong 20 lượt chơi (episode). Kết quả được xem xét qua reward, completion rate, win rate, steps và loss đối với các thuật toán sử dụng mạng nơ-ron. Đề tài chỉ khảo sát một bản đồ và không thực hiện tìm kiếm siêu tham số trên quy mô lớn. Vì vậy, các kết luận được giới hạn trong môi trường và cấu hình thực nghiệm của đề tài, không được xem là kết luận chung về hiệu quả của các thuật toán trong mọi bài toán.

\subsection{Đóng góp của đề tài}

Đóng góp chính của đề tài là xây dựng được môi trường Mini Pacman phục vụ quá trình huấn luyện và đánh giá tác tử RL. Môi trường mô phỏng mê cung, thức ăn, viên sức mạnh, trái cây, ba con ma, số mạng và các điều kiện kết thúc lượt chơi. Trên môi trường này, nhóm cài đặt Q-Learning, DQN và Double DQN trong cùng một ứng dụng, đồng thời xây dựng hệ thống cấu hình YAML, lưu checkpoint, lưu mô hình sau huấn luyện và đánh giá định kỳ.

Từ hệ thống đã xây dựng, nhóm thực hiện chín cấu hình huấn luyện tương ứng với ba thuật toán và ba mức learning rate. Kết quả được tổng hợp qua reward, completion rate, win rate, steps và loss khi phù hợp, qua đó hỗ trợ đánh giá các thuật toán từ nhiều khía cạnh. Bên cạnh các kết quả định lượng, nhóm còn xây dựng biểu đồ so sánh và demo các tác tử sau huấn luyện để quan sát trực tiếp hành vi của mô hình trong môi trường trò chơi.
